\chapter{Local Bindings}
\index{local bindings}
\label{ch:bindings}

In Scheme, local variables are introduced with \emph{binding forms}. Unlike Python, where you create a variable simply by assigning to it, Scheme makes the scope\index{scope} of every variable explicit. This chapter covers \texttt{let}, \texttt{let*}, \texttt{letrec}, internal definitions, and the relationship between \texttt{let} and \texttt{lambda}.

\section{\texttt{let}: Parallel Bindings}
\index{let}

\texttt{let} introduces one or more local variables:

\begin{lstlisting}[style=repl]
> (let ((x 10)
        (y 20))
    (+ x y))
30
\end{lstlisting}

The general form is:

\begin{lstlisting}
(let ((var1 expr1)
      (var2 expr2)
      ...)
  body ...)
\end{lstlisting}

Each \texttt{var} is bound to the value of its \texttt{expr}, and then the \texttt{body} expressions are evaluated in order. The result of the last body expression is the value of the whole \texttt{let}.

\subsection{Bindings Are Parallel}

The key rule: all the \texttt{expr} values are computed \emph{before} any binding takes effect. This means one binding cannot refer to another:

\begin{lstlisting}[style=repl]
> (let ((x 5)
        (y (* x 2)))   ; ERROR: x is not yet bound here
    y)
; Error: unbound variable: x
\end{lstlisting}

This fails because \texttt{x} is not visible to the expression \texttt{(* x 2)}---the bindings are computed in the outer environment. If you need sequential bindings, use \texttt{let*}.

\subsection{Scope Is Limited}

Variables introduced by \texttt{let} exist only within the body. They do not leak into the surrounding scope:

\begin{lstlisting}[style=repl]
> (define x 100)
> (let ((x 5))
    (display x) (newline)
    x)
5
5
> x
100
\end{lstlisting}

The outer \texttt{x} is temporarily shadowed\index{shadowing} inside the \texttt{let}, but remains unchanged afterward.

\subsection{Practical Uses}

Use \texttt{let} to name intermediate results, avoiding repeated computation:

\begin{lstlisting}[style=repl]
> (define (distance x1 y1 x2 y2)
    (let ((dx (- x2 x1))
          (dy (- y2 y1)))
      (sqrt (+ (* dx dx) (* dy dy)))))
> (distance 0 0 3 4)
5
\end{lstlisting}

\section{\texttt{let*}: Sequential Bindings}
\index{let*}

\texttt{let*} is like \texttt{let}, but each binding is visible to subsequent ones. Bindings are evaluated left to right:

\begin{lstlisting}[style=repl]
> (let* ((x 3)
         (y (* x x))
         (z (+ x y)))
    z)
12
\end{lstlisting}

Here, \texttt{y} uses \texttt{x} (which is 3), and \texttt{z} uses both \texttt{x} and \texttt{y}.

\begin{note}[Coming from Python]
\texttt{let*} behaves like consecutive variable assignments in Python:
\begin{Verbatim}[fontsize=\small]
x = 3
y = x * x
z = x + y
\end{Verbatim}
The sequential nature is the same. The difference is that these Scheme bindings are scoped---they vanish after the body, whereas Python variables persist in the enclosing function.
\end{note}

\subsection{When to Use \texttt{let} vs.\ \texttt{let*}}

Use \texttt{let} when bindings are independent of each other. Use \texttt{let*} when later bindings depend on earlier ones. In practice, many programmers default to \texttt{let*} since it is never wrong where \texttt{let} would work, but using \texttt{let} when possible communicates that the bindings are independent.

\section{\texttt{letrec}: Recursive Bindings}
\index{letrec}

\texttt{letrec} allows bindings to refer to each other, enabling mutually recursive local functions:

\begin{lstlisting}[style=repl]
> (letrec ((my-even? (lambda (n)
                       (if (= n 0) #t
                           (my-odd? (- n 1)))))
           (my-odd? (lambda (n)
                      (if (= n 0) #f
                          (my-even? (- n 1))))))
    (list (my-even? 10) (my-odd? 10)))
(#t #f)
\end{lstlisting}

\texttt{my-even?} calls \texttt{my-odd?} and vice versa. Neither \texttt{let} nor \texttt{let*} could express this---\texttt{let} does not allow forward references, and \texttt{let*} only allows references to earlier bindings.

\begin{warning}[Initialization restriction]
The right-hand side expressions in \texttt{letrec} must not evaluate the variables being defined during initialization. In practice, this means the expressions should be \texttt{lambda} forms (which defer evaluation). Do not write \texttt{(letrec ((x (+ y 1)) (y 5)) ...)}---this is undefined behavior.
\end{warning}

\subsection{\texttt{letrec*}}
\index{letrec*}

\texttt{letrec*} is like \texttt{letrec} but guarantees left-to-right evaluation order. It is slightly more permissive: later bindings can use the values of earlier bindings (like \texttt{let*}), and all bindings can be recursive (like \texttt{letrec}):

\begin{lstlisting}[style=repl]
> (letrec* ((x 1)
            (y (+ x 1))
            (f (lambda (n) (+ n x y))))
    (f 10))
13
\end{lstlisting}

\section{Named \texttt{let}: Loops Revisited}
\index{named let}

Named \texttt{let} was introduced in Chapter~\ref{ch:functions}. Here we examine it more carefully as a binding form.

\begin{lstlisting}
(let name ((var init) ...)
  body ...)
\end{lstlisting}

This binds \texttt{name} to a procedure with parameters \texttt{var ...}, calls it with the \texttt{init} values, and makes \texttt{name} available for recursive calls within the body:

\begin{lstlisting}[style=repl]
> (let loop ((numbers '(1 2 3 4 5))
             (sum 0))
    (if (null? numbers)
        sum
        (loop (cdr numbers)
              (+ sum (car numbers)))))
15
\end{lstlisting}

\subsection{Building Lists with Named \texttt{let}}

A common pattern is accumulating results:

\begin{lstlisting}[style=repl]
> (define (squares-up-to n)
    (let loop ((i 1) (acc '()))
      (if (> i n)
          (reverse acc)
          (loop (+ i 1) (cons (* i i) acc)))))
> (squares-up-to 5)
(1 4 9 16 25)
\end{lstlisting}

The accumulator builds the list in reverse (since \texttt{cons} adds to the front), and \texttt{reverse} fixes the order at the end. This is idiomatic Scheme.

\subsection{Named \texttt{let} vs.\ \texttt{do}}

Named \texttt{let} and \texttt{do} can express the same loops. Named \texttt{let} is preferred by most Scheme programmers because the recursive structure is explicit and the control flow is clearer:

\begin{lstlisting}
;; Named let (preferred)
(let loop ((i 0) (sum 0))
  (if (= i 10) sum
      (loop (+ i 1) (+ sum i))))

;; Equivalent do
(do ((i 0 (+ i 1))
     (sum 0 (+ sum i)))
    ((= i 10) sum))
\end{lstlisting}

Both produce 45. Use whichever reads better for your situation.

\section{Internal Definitions}
\index{internal definitions}

Inside a \texttt{lambda}, \texttt{let}, or \texttt{define} body, you can use \texttt{define} to create local bindings:

\begin{lstlisting}[style=repl]
> (define (circle-stats radius)
    (define pi 3.141592653589793)
    (define r-squared (* radius radius))
    (define area (* pi r-squared))
    (define circumference (* 2 pi radius))
    (values area circumference))
> (circle-stats 5)
78.53981633974483
31.41592653589793
\end{lstlisting}

Internal \texttt{define}s are equivalent to \texttt{letrec*}---they can refer to each other and allow mutual recursion. This is the most common way to define local helper functions:

\begin{lstlisting}[style=repl]
> (define (process-list lst)
    (define (helper x)
      (* x x))
    (map helper lst))
> (process-list '(1 2 3 4))
(1 4 9 16)
\end{lstlisting}

\begin{note}[Style preference]
Some programmers prefer internal \texttt{define} for local helpers (reads like top-level code), while others prefer \texttt{let}/\texttt{letrec} (makes scope visually obvious). Both are correct; use what you find clearer.
\end{note}

\section{\texttt{let} Is Syntactic Sugar for \texttt{lambda}}

Understanding the relationship between \texttt{let} and \texttt{lambda} deepens your understanding of both. Every \texttt{let} can be mechanically translated to an immediately-applied \texttt{lambda}:

\begin{lstlisting}
;; These two are identical:
(let ((x 5) (y 10))
  (+ x y))

((lambda (x y) (+ x y)) 5 10)
\end{lstlisting}

The \texttt{let} creates an anonymous function with parameters \texttt{x} and \texttt{y}, then immediately calls it with arguments 5 and 10. This is not just a theoretical curiosity---it explains why \texttt{let} bindings are parallel (arguments to a function call are evaluated independently) and why the body has its own scope.

\section{\texttt{let-values}: Destructuring Multiple Values}
\index{let-values}

When a function returns multiple values (using \texttt{values}), you can bind them with \texttt{let-values}:

\begin{lstlisting}[style=repl]
> (define (min-max lst)
    (let loop ((rest (cdr lst))
               (lo (car lst))
               (hi (car lst)))
      (if (null? rest)
          (values lo hi)
          (loop (cdr rest)
                (min lo (car rest))
                (max hi (car rest))))))
> (let-values (((lo hi) (min-max '(3 1 4 1 5 9 2 6))))
    (display "Min: ") (display lo) (newline)
    (display "Max: ") (display hi) (newline))
Min: 1
Max: 9
\end{lstlisting}

\section{Choosing the Right Binding Form}

\begin{center}
\begin{tabular}{lp{8cm}}
\textbf{Form} & \textbf{When to use} \\
\hline
\texttt{let} & Independent bindings that do not depend on each other \\
\texttt{let*} & Sequential bindings where later ones depend on earlier ones \\
\texttt{letrec} & Mutually recursive local functions \\
\texttt{letrec*} & Mix of sequential values and recursive functions \\
Named \texttt{let} & Tail-recursive loops with named state \\
Internal \texttt{define} & Local helper functions; equivalent to \texttt{letrec*} \\
\texttt{let-values} & Binding multiple return values \\
\end{tabular}
\end{center}

\section*{Summary}

\begin{itemize}
    \item \texttt{let} binds variables in parallel, while \texttt{let*} binds them sequentially.
    \item \texttt{letrec} and \texttt{letrec*} let bindings refer to one another, as for mutually recursive local functions.
    \item Named \texttt{let} is a concise loop, and internal \texttt{define}s behave like \texttt{letrec*}.
    \item \texttt{let} is syntactic sugar for an immediately applied \texttt{lambda}.
    \item \texttt{let-values} binds the results of an expression that returns multiple values.
\end{itemize}

\section*{Exercises}

\begin{exercise}[Scoping Quiz]
Predict the result of each expression below, then check your answers in the REPL. Pay careful attention to which binding form is used:

\begin{lstlisting}
;; Expression A
(let ((x 1))
  (let ((x 2)
        (y x))
    y))

;; Expression B
(let ((x 1))
  (let* ((x 2)
         (y x))
    y))

;; Expression C
(let ((x 10))
  (let ((x (* x 2))
        (y (+ x 3)))
    (list x y)))
\end{lstlisting}

Why do expressions A and B give different results? What value of \code{x} does \code{y} see in each case?
\end{exercise}

\begin{exercise}[Iterative Factorial]
Write a function \code{factorial} using named \code{let} as an iterative loop with an accumulator. Do not use internal \code{define} for a helper---the named \code{let} should be the entire body:

\begin{lstlisting}[style=repl]
> (factorial 0)
1
> (factorial 5)
120
> (factorial 20)
2432902008176640000
\end{lstlisting}

This should be tail-recursive and handle large inputs thanks to Kaappi's bignum support.
\end{exercise}

\begin{exercise}[Euclidean GCD]
Write a function \code{euclidean-gcd} that computes the greatest common divisor of two positive integers using the Euclidean algorithm. Use named \code{let} for the loop:

\begin{lstlisting}[style=repl]
> (euclidean-gcd 48 18)
6
> (euclidean-gcd 100 75)
25
> (euclidean-gcd 17 13)
1
\end{lstlisting}

The algorithm is: if \code{b} is zero, the GCD is \code{a}; otherwise, replace \code{(a, b)} with \code{(b, a mod b)} and repeat. In Python: \texttt{while b: a, b = b, a \% b}. The named \code{let} version maps each ``iteration'' to a tail call.

Kaappi provides a built-in \code{gcd} procedure---the purpose here is to practice named \code{let}.
\end{exercise}

\begin{exercise}[Mutual Recursion with \code{letrec}]
Write a function \code{balanced-parens?} that checks whether a string of parentheses is balanced (e.g., \code{"(()())"} is balanced, \code{"(()"} is not). Use \code{letrec} to define two mutually recursive helpers: \code{expect-open} looks for an opening paren or end of input, and \code{expect-close} tracks the nesting depth:

\begin{lstlisting}[style=repl]
> (balanced-parens? "(()())")
#t
> (balanced-parens? "(()")
#f
> (balanced-parens? "")
#t
> (balanced-parens? ")(")
#f
\end{lstlisting}

Hint: convert the string to a list of characters with \code{string->list}, then walk the list recursively. One approach uses a single recursive function with a depth counter, but the exercise asks for \code{letrec} with two cooperating functions to practice mutual recursion.
\end{exercise}

In the next chapter, we turn to higher-order functions---procedures that take functions as arguments or return them as results.

\chapterend
